\section{Resultados del prototipo}
\subsection{Comparación de métodos}
\begin{table}[H]
\centering
\caption{Resultados derivados de la corrida del proyecto.}
\resizebox{0.98\textwidth}{!}{\begin{tabular}{lrrrrr}
\toprule
Método & Prob. media & Distancia total (NM) & Combustible ref. (L) & Zonas & Costo ref. (S/) \\
\midrule
Greedy & 0.9025 & 1,070.41 & 17,940.01 & 13 & 89,700 \\
MILP   & 0.9352 &   598.31 & 10,027.70 &  6 & 50,138 \\
MARL   & 0.9265 & 2,022.92 & 33,904.11 &  5 & 169,521 \\
\bottomrule
\end{tabular}}
\end{table}
El costo referencial utiliza S/ 5.00 por litro únicamente como escenario ilustrativo.
\begin{figure}[H]\centering
\begin{subfigure}{0.49\textwidth}\includegraphics[width=\linewidth]{figures/probabilidad_metodos.jpg}\caption{Probabilidad media.}\end{subfigure}\hfill
\begin{subfigure}{0.49\textwidth}\includegraphics[width=\linewidth]{figures/distancia_metodos.jpg}\caption{Distancia total.}\end{subfigure}
\vspace{0.3cm}
\begin{subfigure}{0.49\textwidth}\includegraphics[width=\linewidth]{figures/combustible_metodos.jpg}\caption{Combustible de referencia.}\end{subfigure}\hfill
\begin{subfigure}{0.49\textwidth}\includegraphics[width=\linewidth]{figures/costo_combustible_escenario.jpg}\caption{Costo ilustrativo.}\end{subfigure}
\caption{Comparación cuantitativa de Greedy, MILP y MARL.}
\end{figure}
MILP alcanza la mayor probabilidad media y, al mismo tiempo, la menor distancia y combustible de referencia. En un escenario estático, esta ventaja es consistente con el carácter centralizado del optimizador. MARL mantiene una probabilidad alta, pero presenta recorridos mayores; su valor esperado se encuentra en escenarios dinámicos con actualizaciones, cierres o reasignaciones.
\subsection{Asignación espacial MARL}
\begin{figure}[H]\centering
\includegraphics[width=0.86\textwidth]{figures/mapa_final_modelo.jpg}
\caption{Visualización reproducible de la asignación MARL. Las líneas conectan cada puerto con la zona asignada para interpretar la decisión; las rutas operativas deben calcularse con A* y la máscara legal activa.}
\end{figure}
\subsection{Diagnóstico por puerto}
\begin{table}[H]\centering\caption{Resumen MARL por puerto.}
\begin{tabular}{lrrr}\toprule Puerto & Prob. media & Distancia total (NM) & Combustible ref. (L) \\ \midrule
Malabrigo & 0.9145 & 507.98 & 8,513.68 \\
Chimbote & 0.9135 & 477.74 & 8,006.93 \\
Callao & 0.9515 & 1,037.20 & 17,383.50 \\
\bottomrule\end{tabular}\end{table}
\begin{figure}[H]\centering
\begin{subfigure}{0.49\textwidth}\includegraphics[width=\linewidth]{figures/combustible_puerto_marl.jpg}\caption{Combustible por puerto.}\end{subfigure}\hfill
\begin{subfigure}{0.49\textwidth}\includegraphics[width=\linewidth]{figures/ocupacion_marl.jpg}\caption{Ocupación por zona.}\end{subfigure}
\vspace{0.3cm}
\begin{subfigure}{0.62\textwidth}\includegraphics[width=\linewidth]{figures/tradeoff_marl.jpg}\caption{Distancia vs probabilidad.}\end{subfigure}
\caption{Diagnósticos de la política multiagente.}
\end{figure}
\subsection{Entrenamiento}
\begin{figure}[H]\centering\includegraphics[width=0.80\textwidth]{figures/entrenamiento_marl.jpg}\caption{Historial de recompensa del entrenamiento MARL sobre escenarios simulados a partir del mapa real.}\end{figure}
La recompensa se mantiene estable y alcanza rápidamente una meseta. Esto sugiere que el siguiente ciclo de mejora debe revisar la función de recompensa, en particular la penalización por distancia.
\subsection{Simulación}
\begin{figure}[H]\centering\includegraphics[width=0.90\textwidth]{figures/simulacion_montaje.jpg}\caption{Cuatro momentos de una simulación visual del desplazamiento planeado de los 15 agentes. No representa tracking histórico.}\end{figure}